% Simulazione d'esame: prova del 17/07/2023, esercizio 1 (B72-E01) -- ugello aerospike in schlieren.
% Stessi numeri e figure dello svolgimento di riferimento:
% svolgimenti/2023-07-17_B72-E01_aerospike (calcoli.py e figure/).
% Compilare con: bash simulazione_esame/compila.sh 2023-07-17_B72-E01_aerospike

% Preambolo comune delle simulazioni d'esame (vedi GUIDA.md).
% Ogni esercizio: % Preambolo comune delle simulazioni d'esame (vedi GUIDA.md).
% Ogni esercizio: % Preambolo comune delle simulazioni d'esame (vedi GUIDA.md).
% Ogni esercizio: \input{../preambolo} e \graphicspath verso le figure dello
% svolgimento di riferimento; si compila con compila.sh (svolgimento.pdf, A4).

\documentclass[11pt]{article}

\usepackage[T1]{fontenc}
\usepackage[utf8]{inputenc}
\usepackage[italian]{babel}
\usepackage{amsmath,amssymb}
\usepackage{array}
\usepackage{siunitx}
\usepackage{booktabs}
\usepackage{graphicx}
\usepackage{float}
\usepackage{caption}
\usepackage{enumitem}
\usepackage{hyperref}
\usepackage[a4paper,margin=2.5cm]{geometry}

\hypersetup{hidelinks}
\setlength{\parindent}{0pt}
\setlength{\parskip}{4pt}
\allowdisplaybreaks
\sisetup{output-decimal-marker={,}, group-separator={\,}, per-mode=symbol,
         exponent-product=\cdot, range-phrase={\text{--}}}

% Notazione del formulario (come template/svolgimento.tex).
\newcommand{\Cp}{C^{+}}
\newcommand{\Cm}{C^{-}}
\newcommand{\Mr}{M_{r}}
\newcommand{\Ms}{M_{s}}

% Separatore fra relazioni nella stessa formula: \begin{gather*} A \ee B \end{gather*}.
\newcommand{\ee}{,\qquad}

% Titolo dell'esercizio: \esercizio{Prova del GG/MM/AAAA -- Esercizio n}{argomento}.
\newcommand{\esercizio}[2]{%
  \title{#1\\[2pt]\large #2}\author{}\date{}%
  \hypersetup{pdftitle={#1: #2}}}

% Figura centrata: \figura[frazione della larghezza]{file}{didascalia\label{...}}.
\newcommand{\figura}[3][0.7]{%
  \begin{figure}[H]\centering\includegraphics[width=#1\textwidth]{#2}\caption{#3}\end{figure}}

% Testo a sinistra e figura (anche un dettaglio ingrandito) accanto, per tenere la figura vicino al punto che
% si discute: \accanto[larghezza della figura]{file}{didascalia\label{...}}{testo}.
\newcommand{\accanto}[4][0.40]{%
  \par\noindent
  \begin{minipage}[t]{\dimexpr\textwidth-#1\textwidth-0.03\textwidth\relax}\vspace{0pt}#4\end{minipage}\hfill
  \begin{minipage}[t]{#1\textwidth}\vspace{0pt}\centering\includegraphics[width=\linewidth]{#2}%
  \captionof{figure}{#3}\end{minipage}\par\medskip}
 e \graphicspath verso le figure dello
% svolgimento di riferimento; si compila con compila.sh (svolgimento.pdf, A4).

\documentclass[11pt]{article}

\usepackage[T1]{fontenc}
\usepackage[utf8]{inputenc}
\usepackage[italian]{babel}
\usepackage{amsmath,amssymb}
\usepackage{array}
\usepackage{siunitx}
\usepackage{booktabs}
\usepackage{graphicx}
\usepackage{float}
\usepackage{caption}
\usepackage{enumitem}
\usepackage{hyperref}
\usepackage[a4paper,margin=2.5cm]{geometry}

\hypersetup{hidelinks}
\setlength{\parindent}{0pt}
\setlength{\parskip}{4pt}
\allowdisplaybreaks
\sisetup{output-decimal-marker={,}, group-separator={\,}, per-mode=symbol,
         exponent-product=\cdot, range-phrase={\text{--}}}

% Notazione del formulario (come template/svolgimento.tex).
\newcommand{\Cp}{C^{+}}
\newcommand{\Cm}{C^{-}}
\newcommand{\Mr}{M_{r}}
\newcommand{\Ms}{M_{s}}

% Separatore fra relazioni nella stessa formula: \begin{gather*} A \ee B \end{gather*}.
\newcommand{\ee}{,\qquad}

% Titolo dell'esercizio: \esercizio{Prova del GG/MM/AAAA -- Esercizio n}{argomento}.
\newcommand{\esercizio}[2]{%
  \title{#1\\[2pt]\large #2}\author{}\date{}%
  \hypersetup{pdftitle={#1: #2}}}

% Figura centrata: \figura[frazione della larghezza]{file}{didascalia\label{...}}.
\newcommand{\figura}[3][0.7]{%
  \begin{figure}[H]\centering\includegraphics[width=#1\textwidth]{#2}\caption{#3}\end{figure}}

% Testo a sinistra e figura (anche un dettaglio ingrandito) accanto, per tenere la figura vicino al punto che
% si discute: \accanto[larghezza della figura]{file}{didascalia\label{...}}{testo}.
\newcommand{\accanto}[4][0.40]{%
  \par\noindent
  \begin{minipage}[t]{\dimexpr\textwidth-#1\textwidth-0.03\textwidth\relax}\vspace{0pt}#4\end{minipage}\hfill
  \begin{minipage}[t]{#1\textwidth}\vspace{0pt}\centering\includegraphics[width=\linewidth]{#2}%
  \captionof{figure}{#3}\end{minipage}\par\medskip}
 e \graphicspath verso le figure dello
% svolgimento di riferimento; si compila con compila.sh (svolgimento.pdf, A4).

\documentclass[11pt]{article}

\usepackage[T1]{fontenc}
\usepackage[utf8]{inputenc}
\usepackage[italian]{babel}
\usepackage{amsmath,amssymb}
\usepackage{array}
\usepackage{siunitx}
\usepackage{booktabs}
\usepackage{graphicx}
\usepackage{float}
\usepackage{caption}
\usepackage{enumitem}
\usepackage{hyperref}
\usepackage[a4paper,margin=2.5cm]{geometry}

\hypersetup{hidelinks}
\setlength{\parindent}{0pt}
\setlength{\parskip}{4pt}
\allowdisplaybreaks
\sisetup{output-decimal-marker={,}, group-separator={\,}, per-mode=symbol,
         exponent-product=\cdot, range-phrase={\text{--}}}

% Notazione del formulario (come template/svolgimento.tex).
\newcommand{\Cp}{C^{+}}
\newcommand{\Cm}{C^{-}}
\newcommand{\Mr}{M_{r}}
\newcommand{\Ms}{M_{s}}

% Separatore fra relazioni nella stessa formula: \begin{gather*} A \ee B \end{gather*}.
\newcommand{\ee}{,\qquad}

% Titolo dell'esercizio: \esercizio{Prova del GG/MM/AAAA -- Esercizio n}{argomento}.
\newcommand{\esercizio}[2]{%
  \title{#1\\[2pt]\large #2}\author{}\date{}%
  \hypersetup{pdftitle={#1: #2}}}

% Figura centrata: \figura[frazione della larghezza]{file}{didascalia\label{...}}.
\newcommand{\figura}[3][0.7]{%
  \begin{figure}[H]\centering\includegraphics[width=#1\textwidth]{#2}\caption{#3}\end{figure}}

% Testo a sinistra e figura (anche un dettaglio ingrandito) accanto, per tenere la figura vicino al punto che
% si discute: \accanto[larghezza della figura]{file}{didascalia\label{...}}{testo}.
\newcommand{\accanto}[4][0.40]{%
  \par\noindent
  \begin{minipage}[t]{\dimexpr\textwidth-#1\textwidth-0.03\textwidth\relax}\vspace{0pt}#4\end{minipage}\hfill
  \begin{minipage}[t]{#1\textwidth}\vspace{0pt}\centering\includegraphics[width=\linewidth]{#2}%
  \captionof{figure}{#3}\end{minipage}\par\medskip}

\graphicspath{{../../svolgimenti/2023-07-17_B72-E01_aerospike/figure/}}

\esercizio{Prova del 17/07/2023 -- Esercizio 1}{Ugello aerospike piano: lettura della schlieren}

\begin{document}
\maketitle

\emph{Traccia.} La fotografia schlieren mostra un flusso supersonico bidimensionale stazionario per un ugello
aerospike. Il rapporto di espansione (pressione totale / pressione esterna) è pari a 25. Descrivere il campo di flusso
e identificare le caratteristiche del flusso etichettate 1-2-3-4. (10/30)

\section*{Che cosa mostra la figura}
La schlieren rende visibili i gradienti di densità: urti e onde di Mach appaiono come linee nette e sottili, gli
strati di scorrimento come bande sfumate che si allargano. La luminosità dipende dal segno della componente del
gradiente normale al coltello, quindi strutture speculari nelle due metà possono avere tono opposto.

Attraverso una finestra circolare si vede un ugello aerospike piano: due feritoie, in alto e in basso, ciascuna fra un
carenaggio esterno (il blocco nero con la faccia destra verticale) e la spina centrale, che si assottiglia verso destra
e termina con una piccola base piana: la spina è troncata. Il gas arriva dalla camera a sinistra, attraversa le
feritoie convergenti, che al labbro (lo spigolo del carenaggio) sono inclinate verso l'asse, e scorre lungo la spina.
La corrente va da sinistra verso destra, le strutture sono nitide e la traccia dice che il moto è stazionario. Il
campo è simmetrico rispetto all'asse: studio la metà superiore, la inferiore è speculare ($\Cp\leftrightarrow\Cm$).

L'unico dato è $p_0/p_a=25$. Non ho scale: riferisco le lunghezze a $h$, semidistanza fra i due labbri, e le pressioni a
$p_a$. Assumo aria, $\gamma=1{,}4$, perché è una prova in galleria a freddo.

\figura[0.78]{annotata.png}{Schlieren della traccia con le strutture che individuo e gli stati (lettere). Ciano: gola
G al labbro; giallo tratteggiato: raggi del ventaglio di Prandtl--Meyer centrato al labbro, punteggiato: dove ne
cadrebbe la coda; verde: contorno del getto (1); rosso: urto 3; magenta: riflessione 4 sulla spina; arancio: urti di
ricompressione della scia (2); azzurro: onde di Mach nella zona C. I numeri neri sono quelli della traccia.\label{fig:ann}}

Traccia affine nel catalogo: B69-E03 (spina troncata in shadowgraph); per il progetto della spina ideale B59-E02. La
risolvo in modo indipendente.

\section*{Ipotesi e modello, a partire da ciò che vedo}
\begin{itemize}
\item Le pareti della feritoia (carenaggio e spina) arrivano al labbro parallele e subito a valle compaiono linee di
  Mach: la gola è efficace, quindi assumo \textbf{$M=1$ al labbro}, con la direzione delle pareti $\theta_0$.
\item Linee di Mach si vedono ovunque fra la spina e il contorno: tutto il getto è supersonico e uso il
  \textbf{modello 2D stazionario supersonico} (caratteristiche, Prandtl--Meyer, urti obliqui). Non vale nella zona di
  ricircolo dietro la base, che è subsonica, né dentro lo strato di taglio.
\item Il lato esterno del getto è un \textbf{contorno libero}: lì $p=p_a$ e l'aria esterna è in quiete. Sulla spina la
  corrente è tangente alla parete.
\item Fuori dagli urti il moto è isentropico, con $p_0$ uniforme in camera. Gli urti che vedo sono sottili e inclinati
  quasi come onde di Mach, quindi deboli: tratto $p_0$ come costante nel getto e lo controllerò. Ne segue che
  \textbf{sul contorno, dove $p=p_a$, il Mach è lo stesso ovunque}, $M_a$.
\item Le onde di Mach sono rettilinee solo nelle zone uniformi: le direzioni le misuro lì.
\end{itemize}

\section*{Le misure}
Con righello e goniometro, prendendo l'asse come la mezzeria fra i due labbri e la base, misuro (angoli antiorari
dall'asse, metà superiore; la metà inferiore dà gli stessi valori con segno opposto):
\begin{table}[H]
\centering\small
\begin{tabular}{lc}
\toprule
Grandezza & Valore misurato \\
\midrule
pareti della feritoia al labbro (direzione della corrente in gola) & $\theta_0\approx\ang{-45}$ \\
contorno 1, primo tratto (fino a circa $0{,}5\,h$) & $\approx\ang{-4}$ \\
contorno 1, a valle & $\theta_C\approx\ang{-8.5}$ \\
urto 3 & $\approx\ang{-28}$ \\
onde di Mach fra 3 e 1 (zona C) & $\approx\ang{12}$ \\
ultimo tratto della spina & $\approx\ang{-9}$ \\
retta dal labbro allo spigolo della base & $\phi_R\approx\ang{-25}$ \\
urto riflesso 4 & $\approx\ang{12}$ \\
urti di ricompressione della scia & $\approx\pm\ang{18}$ \\
\midrule
punto di riflessione R di 3 sulla spina & $x_R\approx1{,}7\,h$ dal labbro \\
base della spina & $x\approx1{,}95\,h$, altezza $\approx0{,}13\,h$ \\
\bottomrule
\end{tabular}
\end{table}

\section*{Il quadro del campo prima dei conti}
In gola $p^*=0{,}528\,p_0$, molto più della pressione ambiente: il getto che esce dalle feritoie è fortemente
sottoespanso e il resto dell'espansione avviene fuori, a valle del labbro. È il principio dell'aerospike: dal lato
esterno non c'è parete ma il contorno libero a $p=p_a$, quindi il getto si adatta da sé alla pressione ambiente.

Al labbro la corrente gira attorno allo spigolo con un ventaglio di Prandtl--Meyer centrato, di famiglia $\Cm$: espande
fino a $p_a$ e ruota verso la direzione assiale, e le sue onde scendono verso la spina. Questo ventaglio non è una linea
ma una regione: lo riconosco dalle linee di Mach che vi cambiano inclinazione. Una spina ideale è la linea di corrente
dell'onda semplice: la parete ruota punto per punto come la corrente, le onde che vi arrivano non si riflettono e il
getto esce uniforme e assiale. La spina reale invece è troncata e non segue esattamente quella linea di corrente: dove
ruota meno di quanto richiede l'onda semplice manda nel getto espansioni $\Cp$.

Queste espansioni arrivano sul contorno libero, lo piegano verso l'asse e si riflettono, a pressione costante, con
natura opposta: compressioni $\Cm$. Le compressioni convergono e si raccolgono in un urto debole, che attraversa il
getto e si riflette sulla spina come urto, cioè con la stessa natura. Infine, alla troncatura, i getti superiore e
inferiore, inclinati verso l'asse, girano attorno agli spigoli della base (espansioni), racchiudono una zona di
ricircolo a pressione minore di $p_a$ e si incontrano più a valle, dove urti di ricompressione li riportano paralleli
all'asse.

\figura[0.86]{schema.pdf}{Schema della metà superiore in unità di $h$, con la spina e i labbri misurati: gola G
(ciano), raggi del ventaglio al labbro e coda prevista (ocra), espansioni dal lato della spina (blu) che si riflettono
sul contorno come compressioni (rosso tratteggiato) convergenti nell'urto 3, riflessione 4, zona di ricircolo W, urto
di ricompressione 2 e scia sull'asse.\label{fig:schema}}

Fisso gli stati che userò:
\begin{table}[H]
\centering\small
\begin{tabular}{cp{0.44\linewidth}p{0.36\linewidth}}
\toprule
Stato & Dove & Come lo si raggiunge \\
\midrule
a & aria ambiente in quiete, $p=p_a$ & -- \\
G & gola al labbro: $M=1$, $\theta=\theta_0$ & convergente \\
A & a valle del ventaglio, lungo il contorno vicino al labbro: $p=p_a$ & ventaglio $\Cm$ centrato al labbro \\
B & fra la spina e 3: zona a due famiglie; stato rappresentativo davanti a 3 & espansioni $\Cp$ dal lato della spina \\
C & fra 3 e il contorno 1, a valle: $p\simeq p_a$ & urto 3 ($\Cm$) \\
D & dietro la riflessione 4, lungo l'ultimo tratto della spina & urto riflesso 4 ($\Cp$) \\
W & ricircolo dietro la base, subsonico, $p_W<p_a$ & espansione agli spigoli della base \\
F & a valle della scia, corrente assiale & urto di ricompressione ($\Cp$) \\
\bottomrule
\end{tabular}
\caption{Stati della metà superiore. Fra 1 e l'aria ambiente c'è lo strato di taglio.}
\end{table}
La zona B, dove sono attive entrambe le famiglie, si potrebbe risolvere con il metodo delle caratteristiche
(compatibilità lungo $\Cp$ e $\Cm$ più la condizione di parete), ma la soluzione non è analitica e non serve per
riconoscere le strutture: ne uso solo uno stato rappresentativo.

\section*{La gola e il contorno del getto}
In gola
\[
\frac{p^*}{p_0}=\Big(\frac{2}{\gamma+1}\Big)^{\frac{\gamma}{\gamma-1}}=0{,}528
\qquad\Rightarrow\qquad \frac{p_G}{p_a}=\frac{p_0}{p_a}\,\frac{p^*}{p_0}=13{,}2 ,
\]
quindi la sottoespansione è forte. Sul contorno $p=p_a$ e $p_0$ è quello di camera:
\begin{gather*}
\frac{p_0}{p_a}=\big(1+\delta M_a^2\big)^{\frac{\gamma}{\gamma-1}}\ \Rightarrow\ M_a=2{,}75 \ee
\nu_a=\nu(M_a)=\ang{44.6} \ee \mu_a=\arcsin\frac1{M_a}=\ang{21.4}.
\end{gather*}

La struttura \textbf{1} è il \textbf{contorno del getto}, cioè lo strato di taglio (\emph{shear layer}) che separa il
getto dall'aria ambiente in quiete. È una banda sfumata che si allarga verso valle perché vi avviene il mescolamento
(è un contatto, linearmente degenere). Vi vale $p=p_a$, quindi lungo tutto il contorno $M=M_a=2{,}75$; non è uno strato
limite, perché lì non c'è parete. Parte quasi assiale dal labbro e poi si inclina verso l'asse: vedo ora perché.

\section*{Al labbro: il ventaglio e la direzione del getto}
Le onde del ventaglio sono $\Cm$; attraversandole mi muovo lungo una $\Cp$, su cui si conserva $\theta-\nu$. Dalla gola
($\nu=0$, $\theta=\theta_0$) al contorno ($\nu=\nu_a$):
\[
\theta_A-\nu_a=\theta_0-0\qquad\Rightarrow\qquad \theta_A=\theta_0+\nu_a\approx\ang{-0.4}\approx0 .
\]
Le feritoie sono inclinate proprio di $\nu(M_a)$, e controllo subito il confronto: $|\theta_0|\approx\ang{45}$ misurato
contro $\nu(M_a)=\ang{44.6}$. È la condizione di progetto della spina ideale ($|\theta_0|=\nu(M_e)$): il solo ventaglio
porterebbe la corrente a $p_a$ e in direzione assiale, quindi l'ugello lavora vicino al progetto per $p_0/p_a=25$.

L'ultima onda del ventaglio è la $\Cm$ dello stato A:
\[
\theta_{\rm coda}=\theta_A-\mu_a\approx\ang{-21.7},
\]
e raggiungerebbe l'asse a $x=h/\tan\mu_a\approx2{,}56\,h$: la spina reale, lunga $1{,}95\,h$, è troncata a circa tre
quarti della lunghezza ideale.

\section*{Il getto non è un'onda semplice: le espansioni dal lato della spina}
Dentro il getto vedo molte linee sottili oblique: sono onde di Mach deboli della famiglia $\Cp$ (nella metà inferiore
$\Cm$), generate dalle piccole irregolarità della spina. Che si vedano caratteristiche di entrambe le famiglie mi dice
che il campo lungo la spina \textbf{non è un'onda semplice}. Ne trovo la causa confrontando la spina con la linea di
corrente dell'onda semplice. Se il campo fosse il solo ventaglio centrato, la retta dal labbro allo spigolo della base
sarebbe una delle sue onde $\Cm$: $\theta-\mu=\phi_R$, con $\theta=\theta_0+\nu$, cioè
\[
\nu(M)-\mu(M)=\phi_R-\theta_0\quad\Rightarrow\quad M\approx2{,}64,\qquad
\theta_{\rm ideale}=\theta_0+\nu(M)\approx\ang{-2.7}.
\]
L'ultimo tratto della spina è invece inclinato di circa \ang{-9}: la parete ruota circa \ang{6} meno di quanto richiede
l'onda semplice, quindi manda nel getto espansioni $\Cp$ di qualche grado.

Queste espansioni le ritrovo sul contorno, che lascia il labbro quasi nella direzione $\theta_A$ e a valle è inclinato
di $\theta_C\approx\ang{-8.5}$. Quando un'espansione $\Cp$ di ampiezza $\Delta\nu$ arriva sul contorno, attraversandola
mi muovo lungo una $\Cm$, dove $\theta+\nu$ è costante: $\theta$ cala di $\Delta\nu$. La riflessa $\Cm$ deve riportare
$p$ a $p_a$, cioè $\nu$ a $\nu_a$; attraversandola mi muovo lungo una $\Cp$, dove $\theta-\nu$ è costante, e $\theta$
cala di altri $\Delta\nu$. In tutto il contorno ruota di $2\Delta\nu$:
\[
\Delta\nu=\frac{\theta_A-\theta_C}{2}\approx\ang{4.1},
\]
dello stesso ordine dei circa \ang{6} di differenza fra la spina reale e quella ideale: le due misure sono coerenti.

Lo stato B è quello vicino al contorno raggiunto dalle espansioni e non ancora dalle compressioni riflesse:
\begin{gather*}
\nu_B=\nu_a+\Delta\nu\approx\ang{48.7}\ \Rightarrow\ M_B\approx2{,}94 \ee
\theta_B=\theta_A-\Delta\nu\approx\ang{-4.4},\\
\frac{p_B}{p_a}=\frac{p_0}{p_a}\big(1+\delta M_B^2\big)^{-\frac{\gamma}{\gamma-1}}\approx0{,}74 :
\end{gather*}
in B il getto è localmente sovraespanso.

\section*{La struttura 3: l'urto delle compressioni riflesse}
La struttura \textbf{3} è una linea netta e rettilinea che parte dalla zona del labbro e attraversa il getto fino alla
spina. Una prima ipotesi sarebbe la coda del ventaglio, ma questa sarebbe a \ang{-21.7}, mentre 3 è inclinata di circa
\ang{-28}: \textbf{3 non è la coda del ventaglio}. È invece l'\textbf{urto obliquo debole} in cui si raccolgono le
compressioni riflesse dal contorno: queste riportano la corrente a $p_a$ e la ruotano di altri $\Delta\nu$, quindi l'urto
(famiglia $\Cm$) devia la corrente di $\theta=\Delta\nu$ a partire da B. Uso la relazione $\theta$--$\sigma$--$M$, che
viene dalle relazioni di salto applicate alla componente normale ($M_{n1}=M_B\sin\sigma$) con la tangenziale
conservata:
\begin{gather*}
\tan\theta=2\cot\sigma\,\frac{M_B^2\sin^2\sigma-1}{M_B^2(\gamma+\cos2\sigma)+2}
\ \Rightarrow\ \sigma\approx\ang{22.8} \ee M_{n1}=M_B\sin\sigma\approx1{,}14,\\
\frac{p_C}{p_B}=1+\frac{2\gamma}{\gamma+1}\big(M_{n1}^2-1\big)\approx1{,}35
\ \Rightarrow\ p_C\approx p_a \ee M_C\approx2{,}75 \ee \frac{p_{0C}}{p_{0B}}\approx0{,}997 .
\end{gather*}
Controllo subito:
\begin{itemize}
\item $M_{n1}\approx1{,}14>1$ e $\Delta s>0$; la deviazione, \ang{4.1}, è minore di $\theta_{\max}(M_B)=\ang{33.7}$ e
  $\sigma=\ang{22.8}$ è minore dell'angolo del punto di $\theta_{\max}$ (\ang{65}): ramo debole. La perdita di
  pressione totale è dello 0,3\,\%: l'urto è debole e quasi isentropico, e l'ipotesi di $p_0$ costante regge.
\item Dietro 3 $p_C\approx p_a$: la condizione sul contorno è rispettata.
\item Il fronte è inclinato di $\theta_B-\sigma\approx\ang{-27.2}$, contro \ang{-28} misurati, mentre la coda del
  ventaglio sarebbe a \ang{-21.7}: la linea 3 è proprio l'urto delle compressioni riflesse.
\end{itemize}
Nella zona C, fra 3 e il contorno, le onde di Mach $\Cp$ sono rettilinee e inclinate di $\theta_C+\mu_C$. Con le due
inclinazioni misurate:
\[
\mu_C=\phi_{\Cp}-\theta_C\approx\ang{20.5}\quad\Rightarrow\quad M_C=\frac1{\sin\mu_C}\approx2{,}9 ,
\]
coerente con $M_a=2{,}75$ (in C $p\simeq p_a$). Per $M$ elevati $\mu$ varia poco con $M$, quindi questo valore è solo
indicativo.

\section*{La struttura 4: la riflessione sulla spina}
L'urto 3 arriva sulla spina poco a monte della troncatura, a $x_R\approx1{,}7\,h$. La struttura \textbf{4} è la sua
riflessione: su una parete l'urto si riflette con la stessa natura, quindi da lì riparte un urto (famiglia $\Cp$ in
alto, $\Cm$ in basso, dove la traccia lo indica) verso il contorno. Vicino alla parete la riflessione occupa una piccola
zona scura: è l'interazione con lo strato limite della spina.

La deviazione dell'incidente, $\approx\ang{4}$, è molto minore dell'angolo di distacco $\theta_D(M_B)\approx\ang{21}$,
quindi la riflessione è \textbf{regolare}, come si vede (nessun gambetto di Mach). La stimo con la corrente lungo
l'ultimo tratto della spina ($\theta_w=\ang{-9}$, $M\simeq M_B$): l'incidente devia di $-\Delta\nu$ con
$\sigma\approx\ang{22.8}$; il riflesso ($\Cp$) riporta la corrente parallela alla parete con $\sigma_r\approx\ang{24.4}$,
quindi
\[
\theta_w-\Delta\nu+\sigma_r\approx\ang{11},\qquad \frac{p_D}{p_a}\approx1{,}3,\qquad M_D\approx2{,}56 ,
\]
contro circa \ang{12} misurati: verificato.

\section*{La struttura 2: la scia della base troncata}
La struttura \textbf{2} è la scia della base. La corrente gira attorno agli spigoli della base con un'espansione di
Prandtl--Meyer, dietro la base c'è una piccola zona di ricircolo W (``poppa fluida'') e dove i due getti si incontrano
due urti di ricompressione li riallineano all'asse; più a valle si vede la scia lungo l'asse. In W la corrente è
subsonica e il modello 2D supersonico non vale: so solo che $p_W<p_a$, perché la corrente vi entra espandendo attorno
agli spigoli.

Stimo la ricompressione con un solo urto $\Cp$ che riporta all'asse il getto della zona C ($M_a$, $\theta_C$), con
deviazione $|\theta_C|$:
\[
\sigma\approx\ang{28.0},\qquad \text{fronte a }\theta_C+\sigma\approx\ang{19.5},\qquad
\frac{p_F}{p_a}\approx1{,}8,\qquad M_F\approx2{,}36 ,
\]
contro circa \ang{18} misurati. La deviazione, \ang{8.5}, è minore di $\theta_{\max}(M_a)=\ang{32.1}$: ramo debole. La
stima ignora l'espansione agli spigoli della base, che prima abbassa la pressione e aumenta la deviazione da recuperare.

Infine controllo che il getto sia tutto supersonico, come richiede il modello: $M$ va da 1 (G) a 2,94 (B), e
$M_C\approx2{,}75$, $M_D\approx2{,}56$, $M_F\approx2{,}36$. Verificato; resta esclusa la zona W, dove il modello non si
applica.

Non sono strutture del campo i riquadri bianchi dei numeri, il bordo della finestra e il fondo, più chiaro sopra e più
scuro sotto per l'illuminazione non uniforme.

\section*{Il piano $p$--$\theta$}
\figura[1.0]{p_theta.pdf}{Piano $p$--$\theta$ della metà superiore, in scala logaritmica e lineare: ventaglio al labbro
($\theta-\nu=\theta_0$, marrone), espansioni dal lato della spina (blu), polare $\Cm$ di B (urto 3, rosso) e polare
$\Cp$ di C (ricompressione, arancio).\label{fig:pth}}
\figura[0.56]{p_theta_dettaglio.pdf}{Dettaglio del piano $p$--$\theta$ vicino a $p_a$ (scala lineare).\label{fig:pthd}}
Le pressioni vanno da $13{,}2\,p_a$ in gola a $0{,}74\,p_a$ in B, più di un ordine di grandezza, quindi riporto entrambe
le scale. In scala logaritmica si vede tutto il percorso: da G il ventaglio scende lungo $\theta-\nu=\theta_0$ e taglia
la retta $p=p_a$ in A, quasi a $\theta=0$; con una pressione ambiente più bassa la curva proseguirebbe oltre A e il
contorno ruoterebbe verso l'esterno. In scala lineare si vede che quasi tutta la caduta di pressione avviene nel
ventaglio e che le onde successive sono piccole perturbazioni attorno a $p_a$. Il dettaglio mostra il ciclo vicino a
$p_a$: da A le espansioni del lato spina (blu, $\theta+\nu$ costante) portano a B, sotto $p_a$; la polare $\Cm$ di B
(urto 3) torna sulla retta $p=p_a$ in C, con $\theta_C$ uguale all'inclinazione misurata del contorno; la polare $\Cp$
di C, percorsa verso $\theta=0$, dà lo stato F dietro la ricompressione. Nell'immagine annotata (fig.~\ref{fig:ann})
la coda prevista del ventaglio (punteggiato) cade sensibilmente sopra la linea 3, mentre i fronti calcolati di 3, 4 e
della ricompressione coincidono con le linee visibili.

\section*{Risposta}
\begin{table}[H]
\centering\small
\begin{tabular}{cp{0.46\linewidth}p{0.38\linewidth}}
\toprule
& Struttura & Valori \\
\midrule
1 & contorno del getto: strato di taglio fra getto e aria in quiete, $p=p_a$ & $M=M_a=2{,}75$; inclinazione da
$\approx0$ a $\approx\ang{-8.5}$ \\
2 & scia della base troncata: espansioni agli spigoli, ricircolo W ($p_W<p_a$), urti di ricompressione, scia
sull'asse & urti a $\approx\pm\ang{18}$; $p_F\approx1{,}8\,p_a$ \\
3 & urto obliquo debole formato dalle compressioni in cui le espansioni del lato spina si riflettono sul contorno
libero & $\sigma\approx\ang{23}$, $M_{n1}\approx1{,}14$, $p_C/p_B\approx1{,}35$ \\
4 & riflessione regolare di 3 sulla spina, poco a monte della base & riflesso a $\approx\ang{12}$, $p_D\approx1{,}3\,p_a$ \\
\midrule
& gola al labbro & $M=1$, $\theta_0\approx\ang{-45}$, $p^*=13{,}2\,p_a$ \\
& stato adattato sul contorno & $M_a=2{,}75$, $\nu_a=\ang{44.6}$, $\mu_a=\ang{21.4}$ \\
\bottomrule
\end{tabular}
\end{table}
$M_a$, $\nu_a$, $\mu_a$ e $p^*/p_a$ sono esatti per i dati; gli altri valori dipendono dagli angoli misurati e sono
stime.

\section*{Discussione}
\begin{itemize}
\item \textbf{Autoadattamento.} Con le feritoie inclinate di $\nu(M_a)$ l'ugello lavora vicino al progetto: il getto
  esce quasi assiale e nel campo ci sono solo onde deboli ($\Delta\nu\approx\ang{4}$, urto 3 con $M_{n1}\approx1{,}14$).
  È il vantaggio dell'aerospike: il lato esterno è il contorno libero a $p=p_a$, quindi a pressioni ambiente diverse il
  getto cambia forma senza gli urti forti e le separazioni di un ugello a campana fuori progetto.
\item \textbf{Il campo non è un'onda semplice.} Le linee di Mach $\Cp$ che attraversano il getto, il contorno che si
  piega verso l'asse e l'urto 3 dicono la stessa cosa: la spina non assorbe del tutto il ventaglio, perché è troncata e
  il suo profilo ruota meno della linea di corrente dell'onda semplice. Contorno, spina e inclinazione di 3 danno tutti
  espansioni di qualche grado, coerenti fra loro.
\item \textbf{Troncatura.} Tagliare la spina a circa tre quarti della lunghezza ideale fa completare una parte
  dell'espansione nella scia: la base vede una pressione $p_W<p_a$, quindi una resistenza di base, e i due getti si
  ricomprimono con urti. È il prezzo di una spina più corta e leggera.
\item \textbf{Limiti.} La zona di ricircolo è subsonica e richiede un modello diverso; lo strato di taglio e lo strato
  limite della spina sono viscosi (la riflessione 4 è allargata dall'interazione con lo strato limite); le pareti
  laterali della galleria introducono effetti tridimensionali. Le stime dipendono da angoli di pochi gradi: per
  esempio $M_C$ dall'angolo di Mach è solo indicativo.
\item Nella metà inferiore si vedono le stesse strutture speculari (3 e 4 sono indicate lì nella traccia); il tono
  diverso delle stesse strutture nelle due metà è un effetto del coltello, non del campo.
\end{itemize}

\end{document}
